\section{企业 GenAI 的关键趋势}

\subsection{AI 民主化}

\begin{importantbox}{开发者取代 AI 专家}
过去构建 AI 应用需要数据科学家和 ML 专家；如今任何开发者——甚至中学生——都可以用 LLM API 构建 AI 应用。这将 AI 开发者群体扩大了 100 倍以上，是 AI 应用爆发式增长的根本原因。
\end{importantbox}

AI 民主化的背后是两层力量：其一是越来越多的基础模型开放 API，其二是低代码/无代码工具把 prompt 设计、数据清洗与 Ops 流程封装成可视化工作流。这两者共同把 AI 应用从少数 data scientist 转移到大批软件工程师与业务团队。

\subsection{数据需求降低}

\begin{itemize}
    \item 过去企业需要花费数月甚至数年收集标注数据
    \item 大模型已经包含海量预训练知识，fine-tuning 所需数据量大幅减少
    \item 企业可以在数天内启动 AI 项目，而非数年
\end{itemize}

减少数据需求不代表可以跳过数据治理。企业仍需建立「数据目录」「数据质量评分」和「差异化存取」策略，防止模型训练时吸入违规内容。很多企业选择先用公开数据做验证，再用小批高质量数据做微调。

\subsection{模型快速迭代}

\begin{itemize}
    \item 每隔几周就有新模型发布，排行榜不断变化
    \item 企业关注\textbf{平台}而非单个模型——需要能灵活切换模型的基础设施
    \item 18 个月前只有一个顶级模型（GPT-4），如今有多个旗鼓相当的选择
\end{itemize}

频繁迭代也带来版本管理难题。Vertex AI 的特性之一是可以在同一个 Agent pipeline 中定义多个版本（可切换），并记录 A/B 实验效果。这样即使某个版本出现 regress，也能迅速回滚。

\subsection{多模态融合}

\begin{itemize}
    \item 单模态模型正在被多模态模型取代（Gemini 原生支持文本、图像、音频、视频）
    \item 企业对多模态输入（文档中的图表、视频内容分析）需求日增
\end{itemize}

多模态融合意味着要同时处理数据同步、对齐和分辨率问题。Gemini 在语音、图像、文本之间建立统一 embedding space，使得一个输入流可以同时被多个 tokenizers 解析，避免跨模态断层。

\subsection{成本与延迟持续下降}

\begin{itemize}
    \item API 调用成本已下降约 1.5 个数量级
    \item 延迟下降使更多实时应用成为可能
    \item 稀疏 MoE 模型（只激活部分参数）进一步降低推理成本
    \item 有趣发现：运行某些 API 模型比运行等效开源模型更便宜（因省去运维成本）
\end{itemize}

成本下降后，企业更愿意走「算法 + 异构硬件」路线：在高峰期启用低延迟 GPU，低峰期开通 cheaper TPU/CPU，甚至把部分请求 offload 到 regional datacenter。这种做法要求 orchestration 层能感知实际请求特征。

\subsection{本章小结}

企业 GenAI 的核心趋势：AI 民主化、数据门槛降低、模型快速迭代、多模态融合、成本持续下降。

